The largest original Hard-to-Soft contrast primarily reflects changing the target representation rather than evidence that detailed minority identities alone are essential. Uniform is a strong control, and Tie-Uniform removes its canonical tie asymmetry. Neither controls entropy or aggregate class marginals. The cross-evaluation matrix additionally shows that predictor comparisons depend on the accepted image set. An own-ranking advantage should not be promoted to a general claim of improved category choice.

The matched-reference table makes another reporting distinction concrete: changing image eligibility and changing labels both alter F1. Full-vote disagreement avoids silently excluding non-majority images, but remains conditional on the finite crowd panel. The zero-supported-mass retention counts are an adverse result that aggregate selected risk conceals. Contempt, unknown and not-a-face have different meanings; a seven-output model cannot resolve them simply by reporting high MSP.
Tie-Uniform selected risk differs from Uniform by +0.210 pp. The residual own-ranking advantage persists under this tie-preserving control within the conditional interval. Because this adaptive check still uses the same source, three seeds and one model family, it is a representation sensitivity rather than independent replication.
For scale, Uniform lowers selected-checkpoint risk relative to Hard by 2.23 pp, compared with 2.61 pp for Soft. The additional 0.38 pp at 2,620 retained images corresponds to about 9.9 fewer expected disagreements with an empirical annotation draw, not that many corrected images: accepted identities differ and losses are fractional. This arithmetic is not a causal fraction of improvement.
